\chapter{Треугольник редукций: обратная спектральная неоднозначность и коцикл}

Проверяется, может ли корреляционное описание восстановить геометрическое
без наблюдаемых входов, произвольного априорного выбора топологии и свободной
тривиализации контрчленами.

\section{Минимальные объекты треугольника}

\begin{definition}[Минимальный треугольник редукций]
\begin{align}
 R_g&=(K,g,\mathcal E,\nabla;\text{спиновые и граничные данные}),\\
 R_s&=(\mathcal H,D;\Spec D,\text{кратности}),\\
 R_c&=(C_\tau,\tau),\qquad C_\tau=e^{-\tau H},\quad\tau>0.
\end{align}
Корреляционный объект не содержит заранее алгебру координат, отношение
локальности или метки носителя.
\end{definition}

Переход \(T_{gs}\) строит оператор типа Дирака или Лапласа из заданной
геометрии. Переход \(T_{sc}\) сохраняет положительный генератор на опоре:
\begin{equation}H=-\tau^{-1}\log C_\tau.\end{equation}

\section{Точный конечный контрпример}

Связная звезда \(K_{1,4}\) и несвязное объединение
\(C_4\sqcup\{\mathrm{pt}\}\) имеют одинаковый характеристический многочлен
матрицы смежности
\begin{equation}\chi_A(\lambda)=\lambda^3(\lambda^2-4).\end{equation}
После общего положительного сдвига
\begin{equation}H_\star=3I-A_\star,\qquad H_\square=3I-A_\square\end{equation}
получаем один спектр \((1,3,3,3,5)\), а значит,
\begin{equation}
 \Spec e^{-\tau H_\star}=\Spec e^{-\tau H_\square},
 \qquad
 \Tr f(H_\star)=\Tr f(H_\square)
\end{equation}
для любого спектрального профиля \(f\). Однако графы различаются связностью
и последовательностями степеней.

\begin{theorem}[Невосстановимость геометрии по корреляционному спектру]
Спектр, тепловой след и любой скалярный спектральный функционал от
\(C_\tau\) не определяют даже конечную связность. Поэтому на минимальных
объектах канонический морфизм \(T_{cg}\) отсутствует.
\end{theorem}

Полная матрица корреляций приобретает геометрический смысл только после
задания представленной алгебры координат или подалгебры локальности. Для
\(C_\tau=e^{-\tau D^2}\) дополнительно имеется точная неоднозначность
\(D\leftrightarrow-D\), а вырожденные собственные подпространства допускают
другие самосопряжённые квадратные корни.

\begin{nogocriterion}[Запрет восстановления положительным квадратным корнем]
Оператор \(+\sqrt{-\tau^{-1}\log C_\tau}\) нельзя объявлять каноническим
оператором Дирака: он стирает нечётную градуировку и ориентацию и не
проверяет условие первого порядка.
\end{nogocriterion}

\section{Развилка обогащения}

Для слабого объекта \(R_c^{\text{слаб}}=(C_\tau,\tau)\) обратная реконструкция
неединственна. Богатый объект
\begin{equation}
 R_c^{\text{богат}}=(\mathcal A,\mathcal H,C_\tau,J,\gamma;
 \text{данные первого порядка и локальности})
\end{equation}
уже содержит почти всю спектральную и геометрическую структуру. Тогда
возврат становится тождественным забыванием и переименованием, а не выводом.

\begin{proposition}[Дилемма обогащения]
Слабый корреляционный объект недостаточен для \(T_{cg}\), а богатый делает
его тавтологическим и не выводит априорный выбор носителя, вес топологии или
конечные контрчлены.
\end{proposition}

\section{Состояние коцикла и вывод}

Предполагавшийся остаток
\begin{equation}
 \Omega_{gsc}=\Gamma(T_{gs})+\Gamma(T_{sc})+\Gamma(T_{cg})
\end{equation}
не определён в минимальном треугольнике, поскольку \(T_{cg}\) не является
функцией. В богатом треугольнике его нулевое значение тавтологично и не
является новым физическим инвариантом.

\begin{protocol}[Двоичный результат]
\begin{enumerate}
 \item минимальный треугольник редукций: \textbf{не замыкается};
 \item нетривиальный \(T_{cg}\) без свободных коэффициентов:
 \textbf{не существует на текущих данных};
 \item коцикл треугольника: \textbf{не определён};
 \item относительный язык кограниц: \textbf{сохранён};
 \item следующая проверка: \path{version5_boundary_parent_trace_freeze_gate}.
\end{enumerate}
\end{protocol}

Машинный сертификат сохранён в
\path{s2t_v5_reduction_triangle_cocycle_gate_results.json}.